% GENERATED FILE. Edit canonical lessons, then run npm run generate:books.
% canonical-source-hash: fnv1a64:ad69e46bb8965d25
% canonical-lessons: LA-C426-pulverulentus, LA-C426-lutulentus, LA-C426-vestigium, LA-C426-lassus, LA-C426-defatigo

\chapter{Dusty, Muddy, Footprint, Weary, To wear out}
\label{ch:la-dusty-muddy-footprint-weary-to-wear-out}
\hlchaptermodality{426}

\begin{chapteropening}
\textbf{By the end of this chapter:} \emph{I can describe the road in bad weather and how a traveller feels: dusty, muddy, footprints, weary, and worn out.}

Complete the last lesson of chapter 426: I can describe the road in bad weather and how a traveller feels: dusty, muddy, footprints, weary, and worn out.
\end{chapteropening}

\section[pulverulentus, pulverulenta …]{pulverulentus, pulverulenta, pulverulentum — dusty, covered with dust}

\label{lesson:LA-C426-pulverulentus}



\begin{quote}
Before the new one: say the Latin for a wheel rut, then the Latin for to spread.
\end{quote}

\subsection*{You'll want to know: pulverulentus, pulverulenta, pulverulentum}
\textbf{pulverulentus, pulverulenta, pulverulentum} — \textquotedblleft{}dusty, covered with dust\textquotedblright{}.

\begin{quote}
Aestāte via pulverulenta est, et viātōrēs aquam quaerunt.
\end{quote}

\begin{quote}
\emph{In summer the road is dusty, and travellers look for water.}
\end{quote}

An adjective of three endings, from \emph{pulvis, pulveris} 'dust' with the ending \emph{-ulentus} 'full of'. English \emph{pulverize} is a cousin.

\begin{cousinweb}
\emph{pulverulentus} is built on \emph{pulvis, pulveris} 'dust' with the suffix \emph{-ulentus} 'full of' (EtymDB; Lewis \& Short). English \emph{pulverize} comes from the same noun.
\end{cousinweb}

\begin{grammarlens}[title={What you've built}]
One more word for telling, explaining and describing at B1.
\end{grammarlens}

\subsection*{Your turn}
\begin{itemize}
\raggedright
  \item \emph{Say it:} \emph{pulverulentus, pulverulenta, pulverulentum}
  \item \emph{Say it:} \emph{pulverulentus, pulverulenta, pulverulentum}, once more
  \item \emph{Say it:} a sentence of your own with \emph{pulverulentus, pulverulenta, pulverulentum}
  \item \emph{Recall:} say the Latin for a wheel rut, then the Latin for to spread, then say \emph{pulverulentus, pulverulenta, pulverulentum} again
\end{itemize}

\begin{tcolorbox}[breakable,colback=teal!4,colframe=teal!35!black,title={Before you move on}]
What does pulverulentus, pulverulenta, pulverulentum mean? (\textquotedblleft{}Dusty, covered with dust\textquotedblright{}.) Say it once more.
\end{tcolorbox}
\section[lutulentus, lutulenta, lutulentum]{lutulentus, lutulenta, lutulentum — muddy, covered with mud}

\label{lesson:LA-C426-lutulentus}



\begin{quote}
Before the new one: say the Latin for to spread, then the Latin for dusty.
\end{quote}

\subsection*{You'll want to know: lutulentus, lutulenta, lutulentum}
\textbf{lutulentus, lutulenta, lutulentum} — \textquotedblleft{}muddy, covered with mud\textquotedblright{}.

\begin{quote}
Post pluviam via lutulenta est, et plaustra lentē eunt.
\end{quote}

\begin{quote}
\emph{After the rain the road is muddy, and the wagons go slowly.}
\end{quote}

An adjective of three endings, from \emph{lutum} 'mud'. Also of a dirty, low person.

\begin{cousinweb}
\emph{lutulentus} is built on \emph{lutum} 'mud' with the suffix \emph{-ulentus} 'full of' (EtymDB; Lewis \& Short), the same pattern as \emph{pulverulentus} 'dusty'.
\end{cousinweb}

\begin{grammarlens}[title={What you've built}]
One more word for telling, explaining and describing at B1.
\end{grammarlens}

\subsection*{Your turn}
\begin{itemize}
\raggedright
  \item \emph{Say it:} \emph{lutulentus, lutulenta, lutulentum}
  \item \emph{Say it:} \emph{lutulentus, lutulenta, lutulentum}, once more
  \item \emph{Say it:} a sentence of your own with \emph{lutulentus, lutulenta, lutulentum}
  \item \emph{Recall:} say the Latin for to spread, then the Latin for dusty, then say \emph{lutulentus, lutulenta, lutulentum} again
\end{itemize}

\begin{tcolorbox}[breakable,colback=teal!4,colframe=teal!35!black,title={Before you move on}]
What does lutulentus, lutulenta, lutulentum mean? (\textquotedblleft{}Muddy, covered with mud\textquotedblright{}.) Say it once more.
\end{tcolorbox}
\section[vestīgium, vestīgiī]{vestīgium, vestīgiī — a footprint, a track; a trace}

\label{lesson:LA-C426-vestigium}



\begin{quote}
Before the new one: say the Latin for dusty, then the Latin for muddy.
\end{quote}

\subsection*{You'll want to know: vestīgium, vestīgiī}
\textbf{vestīgium, vestīgiī} — \textquotedblleft{}a footprint, a track; a trace\textquotedblright{}.

\begin{quote}
In lutō vestīgia equōrum et mūlōrum sunt.
\end{quote}

\begin{quote}
\emph{In the mud there are the footprints of horses and mules.}
\end{quote}

Second declension, neuter: \emph{vestīgium, vestīgiī n.} \emph{Eōdem vestīgiō}, on the same spot. English \emph{vestige} comes from it.

\begin{cousinweb}
The origin of \emph{vestīgium} is not settled: Lewis \& Short call the etymology doubtful, and EtymDB gives no parent. English \emph{vestige} comes from it, and \emph{investigate} from the related verb \emph{investīgāre}.
\end{cousinweb}

\begin{grammarlens}[title={What you've built}]
One more word for telling, explaining and describing at B1.
\end{grammarlens}

\subsection*{Your turn}
\begin{itemize}
\raggedright
  \item \emph{Say it:} \emph{vestīgium, vestīgiī}
  \item \emph{Say it:} \emph{vestīgium, vestīgiī}, once more
  \item \emph{Say it:} a sentence of your own with \emph{vestīgium, vestīgiī}
  \item \emph{Recall:} say the Latin for dusty, then the Latin for muddy, then say \emph{vestīgium, vestīgiī} again
\end{itemize}

\begin{tcolorbox}[breakable,colback=teal!4,colframe=teal!35!black,title={Before you move on}]
What does vestīgium, vestīgiī mean? (\textquotedblleft{}A footprint, a track; a trace\textquotedblright{}.) Say it once more.
\end{tcolorbox}
\section[lassus, lassa, lassum]{lassus, lassa, lassum — weary, tired out}

\label{lesson:LA-C426-lassus}



\begin{quote}
Before the new one: say the Latin for muddy, then the Latin for a footprint.
\end{quote}

\subsection*{You'll want to know: lassus, lassa, lassum}
\textbf{lassus, lassa, lassum} — \textquotedblleft{}weary, tired out\textquotedblright{}.

\begin{quote}
Post longum iter viātor lassus in caupōnā sedet.
\end{quote}

\begin{quote}
\emph{After a long journey the weary traveller sits in the inn.}
\end{quote}

An adjective of three endings, close to \emph{fessus}. English \emph{lassitude} comes from the noun \emph{lassitūdō}.

\begin{cousinweb}
\emph{lassus} is inherited and is usually traced to PIE \emph{led-} 'slow, weary' (EtymDB); Lewis \& Short call its origin doubtful. English \emph{lassitude} comes from the Latin noun \emph{lassitūdō}.
\end{cousinweb}

\begin{grammarlens}[title={What you've built}]
One more word for telling, explaining and describing at B1.
\end{grammarlens}

\subsection*{Your turn}
\begin{itemize}
\raggedright
  \item \emph{Say it:} \emph{lassus, lassa, lassum}
  \item \emph{Say it:} \emph{lassus, lassa, lassum}, once more
  \item \emph{Say it:} a sentence of your own with \emph{lassus, lassa, lassum}
  \item \emph{Recall:} say the Latin for muddy, then the Latin for a footprint, then say \emph{lassus, lassa, lassum} again
\end{itemize}

\begin{tcolorbox}[breakable,colback=teal!4,colframe=teal!35!black,title={Before you move on}]
What does lassus, lassa, lassum mean? (\textquotedblleft{}Weary, tired out\textquotedblright{}.) Say it once more.
\end{tcolorbox}
\section[dēfatīgō, dēfatīgāre]{dēfatīgō, dēfatīgāre — to wear out, to tire out}

\label{lesson:LA-C426-defatigo}



\begin{quote}
Before the new one: say the Latin for a footprint, then the Latin for weary.
\end{quote}

\subsection*{You'll want to know: dēfatīgō, dēfatīgāre}
\textbf{dēfatīgō, dēfatīgāre} — \textquotedblleft{}to wear out, to tire out\textquotedblright{}.

\begin{quote}
Via longa equōs dēfatīgat.
\end{quote}

\begin{quote}
\emph{The long road wears the horses out.}
\end{quote}

\emph{dēfatīgō, dēfatīgāre, dēfatīgāvī, dēfatīgātum}. Often passive: \emph{dēfatīgātus} 'worn out'. The simple verb \emph{fatīgō} gives English \emph{fatigue}.

\begin{cousinweb}
\emph{dēfatīgō} is built from \emph{dē-} 'down, completely' and \emph{fatīgō} 'to tire, to weary' (EtymDB; Lewis \& Short). English \emph{fatigue} comes, through French, from \emph{fatīgō}.
\end{cousinweb}

\begin{grammarlens}[title={What you've built}]
One more word for telling, explaining and describing at B1.
\end{grammarlens}

\subsection*{Your turn}
\begin{itemize}
\raggedright
  \item \emph{Say it:} \emph{dēfatīgō, dēfatīgāre}
  \item \emph{Say it:} \emph{dēfatīgō, dēfatīgāre}, once more
  \item \emph{Say it:} a sentence of your own with \emph{dēfatīgō, dēfatīgāre}
  \item \emph{Recall:} say the Latin for a footprint, then the Latin for weary, then say \emph{dēfatīgō, dēfatīgāre} again
\end{itemize}

\begin{tcolorbox}[breakable,colback=teal!4,colframe=teal!35!black,title={Before you move on}]
What does dēfatīgō, dēfatīgāre mean? (\textquotedblleft{}To wear out, to tire out\textquotedblright{}.) Say it once more.
\end{tcolorbox}
